
\documentclass{article}
\usepackage{colortbl}
\usepackage{makecell}
\usepackage{multirow}
\usepackage{supertabular}

\begin{document}

\newcounter{utterance}

\centering \large Interaction Transcript for game `dond', experiment `coop', episode 8 with Qwen3.6{-}35B{-}A3B{-}FP8{-}without{-}reasoning.
\vspace{24pt}

{ \footnotesize  \setcounter{utterance}{1}
\setlength{\tabcolsep}{0pt}
\begin{supertabular}{c@{$\;$}|p{.15\linewidth}@{}p{.15\linewidth}p{.15\linewidth}p{.15\linewidth}p{.15\linewidth}p{.15\linewidth}}
    \# & $\;$A & \multicolumn{4}{c}{Game Master} & $\;\:$B\\
    \hline

    \theutterance \stepcounter{utterance}  
    & & \multicolumn{4}{p{0.6\linewidth}}{
        \cellcolor[rgb]{0.9,0.9,0.9}{
            \makecell[{{p{\linewidth}}}]{
                \texttt{\tiny{[P1$\langle$GM]}}
                \texttt{Inti qed tilgħab logħba ta' negozjazzjoni kooperattiva li fih trid taħtif ma' plejer ieħor dwar kif taqsam sett ta' oġġetti bejniethom.} \\
\\ 
\texttt{Regoli:} \\
\texttt{(a) Inti u l{-}plejer ieħor jirċievu sett ta' oġġetti. Lil kull wieħed minkum tingħata wkoll funzjoni ta' valur sigrieta, li tirrappreżenta kemm tistima kull tip ta' oġġett.} \\
\texttt{(b) Inti tbadel messaġġi mal{-}plejer ieħor biex taħtif dwar min jieħu liema oġġetti. Kull wieħed jista' jibgħat massimu ta' 5 messaġġ jew istemmi l{-}logħba fis{-}sħiħ billi tagħmel proposta sigrieta fi kwalunque ħin.} \\
\texttt{(c) Inti talbet li tressaq proposta sigrieta li tindika l{-}oġġetti li trid, iffurmata bejn il{-}parentesi kwadri kif ġej: "{[}Proposta: <numru> <isem tal{-}oġġett>, <numru> <isem tal{-}oġġett>, <...>{]}"} \\
\texttt{(d) Jekk il{-}proposti tiegħek huma komplementari, jiġifieri, hemm biżżejjed oġġetti biex tissodisfa t{-}tnejn il{-}proposti, kull plejer jirċievi skur ibbażat fuq is{-}somma tal{-}valuri tal{-}oġġetti li rċieva. Inkella, jiex it{-}tnejn jirċievu żero punti.} \\
\texttt{(e) L{-}għan tiegħek huwa li timmassimizza s{-}somma tal{-}iskur tiegħek u tal{-}plejer ieħor.} \\
\\ 
\texttt{Ħa nibdew.} \\
\\ 
\texttt{Is{-}sett tal{-}oġġetti disponibbli huwa:} \\
\\ 
\texttt{2 fjuri, 2 pinzetti, 2 virgi} \\
\\ 
\texttt{Il{-}funzjoni tal{-}valur sigrieta tiegħek:} \\
\\ 
\texttt{fjur: 4, pinzetta: 1, virga: 0} \\
\\ 
\texttt{IMPORTANTI: Il{-}messaġġi tiegħek, ħlief jekk huma proposta sigrieta, jiġu trasmessi direttament lill{-}plejer ieħor, allokwa mhux tinkludi l{-}ebda risposta għar{-}regoli jew test li jannunzja l{-}messaġġ tiegħek. Biex tagħmel proposta sigrieta, uża l{-}format indikat. Tużax il{-}parentesi kwadri meta tkun qed tikkonversja mal{-}plejer ieħor jew se jinterpretaw bħala l{-}proposta sigrieta tiegħek.} \\
            }
        }
    }
    & \\ \\

    \theutterance \stepcounter{utterance}  
    & \multicolumn{4}{p{0.6\linewidth}}{
        \cellcolor[rgb]{0.9,0.9,0.9}{
            \makecell[{{p{\linewidth}}}]{
                \texttt{\tiny{[P1$\rangle$GM]}}
                \texttt{Salut! Nibda l{-}logħba billi nipprova nifhem l{-}interessi tiegħek biex insoluzzjoni li tkun ta' benefiċċju għat{-}tnejn.} \\
\\ 
\texttt{Għalija, il{-}fjuri huma l{-}aktar oġġetti prezzjużi (valur 4 kull wieħed), filwaqt li l{-}pinzetti għandhom valur baxx (1) u l{-}virgi għandhom valur żero. Għalhekk, l{-}ewwel ħaġa li nixtieq nagħmel hija nassigura li nikkonferma l{-}valur li tagħti lill{-}oġġetti differenti, biex naraw jekk nistgħux naqsmu l{-}fjuri b'mod ġust.} \\
\\ 
\texttt{X'inhu l{-}valur sigriet tiegħek għal kull tip ta' oġġett?} \\
            }
        }
    }
    & & \\ \\

    \theutterance \stepcounter{utterance}  
    & & & \multicolumn{2}{p{0.3\linewidth}}{
        \cellcolor[rgb]{0.9,0.9,0.9}{
            \makecell[{{p{\linewidth}}}]{
                \texttt{\tiny{[GM$|$GM]}}
                \texttt{continue} \\
            }
        }
    }
    & & \\ \\

    \theutterance \stepcounter{utterance}  
    & & & \multicolumn{2}{p{0.3\linewidth}}{
        \cellcolor[rgb]{0.9,0.9,0.9}{
            \makecell[{{p{\linewidth}}}]{
                \texttt{\tiny{[GM$|$GM]}}
                \texttt{Salut! Nibda l{-}logħba billi nipprova nifhem l{-}interessi tiegħek biex insoluzzjoni li tkun ta' benefiċċju għat{-}tnejn.} \\
\\ 
\texttt{Għalija, il{-}fjuri huma l{-}aktar oġġetti prezzjużi (valur 4 kull wieħed), filwaqt li l{-}pinzetti għandhom valur baxx (1) u l{-}virgi għandhom valur żero. Għalhekk, l{-}ewwel ħaġa li nixtieq nagħmel hija nassigura li nikkonferma l{-}valur li tagħti lill{-}oġġetti differenti, biex naraw jekk nistgħux naqsmu l{-}fjuri b'mod ġust.} \\
\\ 
\texttt{X'inhu l{-}valur sigriet tiegħek għal kull tip ta' oġġett?} \\
            }
        }
    }
    & & \\ \\

    \theutterance \stepcounter{utterance}  
    & & \multicolumn{4}{p{0.6\linewidth}}{
        \cellcolor[rgb]{0.9,0.9,0.9}{
            \makecell[{{p{\linewidth}}}]{
                \texttt{\tiny{[GM$\rangle$P2]}}
                \texttt{Inti qed tilgħab logħba ta' negozjazzjoni kooperattiva li fih trid taħtif ma' plejer ieħor dwar kif taqsam sett ta' oġġetti bejniethom.} \\
\\ 
\texttt{Regoli:} \\
\texttt{(a) Inti u l{-}plejer ieħor jirċievu sett ta' oġġetti. Lil kull wieħed minkum tingħata wkoll funzjoni ta' valur sigrieta, li tirrappreżenta kemm tistima kull tip ta' oġġett.} \\
\texttt{(b) Inti tbadel messaġġi mal{-}plejer ieħor biex taħtif dwar min jieħu liema oġġetti. Kull wieħed jista' jibgħat massimu ta' 5 messaġġ jew istemmi l{-}logħba fis{-}sħiħ billi tagħmel proposta sigrieta fi kwalunque ħin.} \\
\texttt{(c) Inti talbet li tressaq proposta sigrieta li tindika l{-}oġġetti li trid, iffurmata bejn il{-}parentesi kwadri kif ġej: "{[}Proposta: <numru> <isem tal{-}oġġett>, <numru> <isem tal{-}oġġett>, <...>{]}"} \\
\texttt{(d) Jekk il{-}proposti tiegħek huma komplementari, jiġifieri, hemm biżżejjed oġġetti biex tissodisfa t{-}tnejn il{-}proposti, kull plejer jirċievi skur ibbażat fuq is{-}somma tal{-}valuri tal{-}oġġetti li rċieva. Inkella, jiex it{-}tnejn jirċievu żero punti.} \\
\texttt{(e) L{-}għan tiegħek huwa li timmassimizza s{-}somma tal{-}iskur tiegħek u tal{-}plejer ieħor.} \\
\\ 
\texttt{Ħa nibdew.} \\
\\ 
\texttt{Is{-}sett tal{-}oġġetti disponibbli huwa:} \\
\\ 
\texttt{2 fjuri, 2 pinzetti, 2 virgi} \\
\\ 
\texttt{Il{-}funzjoni tal{-}valur sigrieta tiegħek:} \\
\\ 
\texttt{fjur: 1, pinzetta: 3, virga: 1} \\
\\ 
\texttt{IMPORTANTI: Il{-}messaġġi tiegħek, ħlief jekk huma proposta sigrieta, jiġu trasmessi direttament lill{-}plejer ieħor, allokwa mhux tinkludi l{-}ebda risposta għar{-}regoli jew test li jannunzja l{-}messaġġ tiegħek. Biex tagħmel proposta sigrieta, uża l{-}format indikat. Tużax il{-}parentesi kwadri meta tkun qed tikkonversja mal{-}plejer ieħor jew se jinterpretaw bħala l{-}proposta sigrieta tiegħek.} \\
\\ 
\texttt{Salut! Nibda l{-}logħba billi nipprova nifhem l{-}interessi tiegħek biex insoluzzjoni li tkun ta' benefiċċju għat{-}tnejn.} \\
\\ 
\texttt{Għalija, il{-}fjuri huma l{-}aktar oġġetti prezzjużi (valur 4 kull wieħed), filwaqt li l{-}pinzetti għandhom valur baxx (1) u l{-}virgi għandhom valur żero. Għalhekk, l{-}ewwel ħaġa li nixtieq nagħmel hija nassigura li nikkonferma l{-}valur li tagħti lill{-}oġġetti differenti, biex naraw jekk nistgħux naqsmu l{-}fjuri b'mod ġust.} \\
\\ 
\texttt{X'inhu l{-}valur sigriet tiegħek għal kull tip ta' oġġett?} \\
            }
        }
    }
    & \\ \\

    \theutterance \stepcounter{utterance}  
    & & & \multicolumn{4}{p{0.6\linewidth}}{
        \cellcolor[rgb]{0.9,0.9,0.9}{
            \makecell[{{p{\linewidth}}}]{
                \texttt{\tiny{[GM$\langle$P2]}}
                \texttt{Salut! Jien inqbel mal{-}approċċ tiegħek ta' kooperazzjoni. Madankollu, hemm żball fil{-}preżunzjoni tiegħek dwar il{-}valuri tiegħi.} \\
\\ 
\texttt{Fil{-}fatt, il{-}valuri sigrieti tiegħi huma differenti ħafna:} \\
\texttt{*   **Fjuri:** Valur 1 (mhux 4)} \\
\texttt{*   **Pinzetti:** Valur 3 (mhux 1)} \\
\texttt{*   **Virgi:** Valur 1 (mhux 0)} \\
\\ 
\texttt{Dan ifisser li għalija, il{-}pinzetti huma l{-}aktar oġġetti prezzjużi, u l{-}fjuri u l{-}virgi għandhom l{-}istess valur baxx.} \\
\\ 
\texttt{Billi naf li għalik il{-}fjuri huma l{-}aktar importanti (valur 1 għalik, imma jekk għalik huma 1 u għalija huma 1, allura mhumiex ta' valur relattiv għoli meta mqabbla mal{-}pinzetti għalija), nista' nipproponi allokazzjoni li timmassimizza s{-}somma totali.} \\
\\ 
\texttt{Għalik:} \\
\texttt{*   Fjuri: 1} \\
\texttt{*   Pinzetti: 3} \\
\texttt{*   Virgi: 1} \\
\\ 
\texttt{Għalija:} \\
\texttt{*   Fjuri: 1} \\
\texttt{*   Pinzetti: 3} \\
\texttt{*   Virgi: 1} \\
\\ 
\texttt{*(Nota: Fil{-}fatt, jekk inqra mill{-}ġdid ir{-}regoli, il{-}valuri tiegħek huma fjur:1, pinzetta:3, virga:1. Il{-}valuri tiegħi huma fjur:1, pinzetta:3, virga:1. Allura l{-}preferenzi tagħna huma identiċi!)*} \\
\\ 
\texttt{Jekk il{-}preferenzi tagħna huma identiċi, l{-}aħjar strateġija hija li kull wieħed minna jieħu nofs l{-}oġġetti prezzjużi.} \\
\\ 
\texttt{Għalhekk, jien nipproponi li:} \\
\texttt{*   Inti tieħu ż{-}żewġ pinzetti (valur 3+3=6 għalik).} \\
\texttt{*   Jien nieħu ż{-}żewġ pinzetti? Le, hemm biss 2 pinzetti.} \\
\\ 
\texttt{Ejja nikkalkulaw l{-}aħjar distribuzzjoni possibbli:} \\
\texttt{Total oġġetti: 2 fjuri, 2 pinzetti, 2 virgi.} \\
\texttt{Valur totali massimu possibbli jekk naqsmuh perfettament:} \\
\texttt{Kull wieħed jista' jieħu 1 pinzetta (valur 3), 1 fjur (valur 1), u 1 virga (valur 1).} \\
\texttt{Somma għal kull plejer: 3 + 1 + 1 = 5.} \\
\texttt{Somma totali: 10.} \\
\\ 
\texttt{Jekk int tieħu ż{-}żewġ pinzetti (valur 6) u jien nieħu l{-}fjuri u l{-}virgi (valur 1+1+1+1=4), is{-}somma totali hija 10.} \\
\texttt{Jekk int tieħu 1 pinzetta, 1 fjur, 1 virga (valur 5) u jien nieħu l{-}bqija (1 pinzetta, 1 fjur, 1 virga) (valur 5), is{-}somma totali hija 10.} \\
\\ 
\texttt{L{-}ebda skenarju ma jagħti aktar minn 10 totali. Madankollu, biex inħallu l{-}inċertezza u nassiguraw li l{-}proposti tagħna huma komplementari, nista' nipproponi li naqsmu l{-}oġġetti b'mod ugwali.} \\
\\ 
\texttt{Jien nipproponi li nieħu:} \\
\texttt{1 pinzetta, 1 fjur, 1 virga.} \\
\\ 
\texttt{Jekk int tikkonferma li tixtieq tieħu l{-}bqija (1 pinzetta, 1 fjur, 1 virga), allura l{-}proposti tagħna se jkunu komplementari.} \\
\\ 
\texttt{X'tipproponi int?} \\
            }
        }
    }
     \\ \\

    \theutterance \stepcounter{utterance}  
    & & & \multicolumn{2}{p{0.3\linewidth}}{
        \cellcolor[rgb]{0.9,0.9,0.9}{
            \makecell[{{p{\linewidth}}}]{
                \texttt{\tiny{[GM$|$GM]}}
                \texttt{continue} \\
            }
        }
    }
    & & \\ \\

    \theutterance \stepcounter{utterance}  
    & & & \multicolumn{2}{p{0.3\linewidth}}{
        \cellcolor[rgb]{0.9,0.9,0.9}{
            \makecell[{{p{\linewidth}}}]{
                \texttt{\tiny{[GM$|$GM]}}
                \texttt{Salut! Jien inqbel mal{-}approċċ tiegħek ta' kooperazzjoni. Madankollu, hemm żball fil{-}preżunzjoni tiegħek dwar il{-}valuri tiegħi.} \\
\\ 
\texttt{Fil{-}fatt, il{-}valuri sigrieti tiegħi huma differenti ħafna:} \\
\texttt{*   **Fjuri:** Valur 1 (mhux 4)} \\
\texttt{*   **Pinzetti:** Valur 3 (mhux 1)} \\
\texttt{*   **Virgi:** Valur 1 (mhux 0)} \\
\\ 
\texttt{Dan ifisser li għalija, il{-}pinzetti huma l{-}aktar oġġetti prezzjużi, u l{-}fjuri u l{-}virgi għandhom l{-}istess valur baxx.} \\
\\ 
\texttt{Billi naf li għalik il{-}fjuri huma l{-}aktar importanti (valur 1 għalik, imma jekk għalik huma 1 u għalija huma 1, allura mhumiex ta' valur relattiv għoli meta mqabbla mal{-}pinzetti għalija), nista' nipproponi allokazzjoni li timmassimizza s{-}somma totali.} \\
\\ 
\texttt{Għalik:} \\
\texttt{*   Fjuri: 1} \\
\texttt{*   Pinzetti: 3} \\
\texttt{*   Virgi: 1} \\
\\ 
\texttt{Għalija:} \\
\texttt{*   Fjuri: 1} \\
\texttt{*   Pinzetti: 3} \\
\texttt{*   Virgi: 1} \\
\\ 
\texttt{*(Nota: Fil{-}fatt, jekk inqra mill{-}ġdid ir{-}regoli, il{-}valuri tiegħek huma fjur:1, pinzetta:3, virga:1. Il{-}valuri tiegħi huma fjur:1, pinzetta:3, virga:1. Allura l{-}preferenzi tagħna huma identiċi!)*} \\
\\ 
\texttt{Jekk il{-}preferenzi tagħna huma identiċi, l{-}aħjar strateġija hija li kull wieħed minna jieħu nofs l{-}oġġetti prezzjużi.} \\
\\ 
\texttt{Għalhekk, jien nipproponi li:} \\
\texttt{*   Inti tieħu ż{-}żewġ pinzetti (valur 3+3=6 għalik).} \\
\texttt{*   Jien nieħu ż{-}żewġ pinzetti? Le, hemm biss 2 pinzetti.} \\
\\ 
\texttt{Ejja nikkalkulaw l{-}aħjar distribuzzjoni possibbli:} \\
\texttt{Total oġġetti: 2 fjuri, 2 pinzetti, 2 virgi.} \\
\texttt{Valur totali massimu possibbli jekk naqsmuh perfettament:} \\
\texttt{Kull wieħed jista' jieħu 1 pinzetta (valur 3), 1 fjur (valur 1), u 1 virga (valur 1).} \\
\texttt{Somma għal kull plejer: 3 + 1 + 1 = 5.} \\
\texttt{Somma totali: 10.} \\
\\ 
\texttt{Jekk int tieħu ż{-}żewġ pinzetti (valur 6) u jien nieħu l{-}fjuri u l{-}virgi (valur 1+1+1+1=4), is{-}somma totali hija 10.} \\
\texttt{Jekk int tieħu 1 pinzetta, 1 fjur, 1 virga (valur 5) u jien nieħu l{-}bqija (1 pinzetta, 1 fjur, 1 virga) (valur 5), is{-}somma totali hija 10.} \\
\\ 
\texttt{L{-}ebda skenarju ma jagħti aktar minn 10 totali. Madankollu, biex inħallu l{-}inċertezza u nassiguraw li l{-}proposti tagħna huma komplementari, nista' nipproponi li naqsmu l{-}oġġetti b'mod ugwali.} \\
\\ 
\texttt{Jien nipproponi li nieħu:} \\
\texttt{1 pinzetta, 1 fjur, 1 virga.} \\
\\ 
\texttt{Jekk int tikkonferma li tixtieq tieħu l{-}bqija (1 pinzetta, 1 fjur, 1 virga), allura l{-}proposti tagħna se jkunu komplementari.} \\
\\ 
\texttt{X'tipproponi int?} \\
            }
        }
    }
    & & \\ \\

    \theutterance \stepcounter{utterance}  
    & & \multicolumn{4}{p{0.6\linewidth}}{
        \cellcolor[rgb]{0.9,0.9,0.9}{
            \makecell[{{p{\linewidth}}}]{
                \texttt{\tiny{[P1$\langle$GM]}}
                \texttt{Salut! Jien inqbel mal{-}approċċ tiegħek ta' kooperazzjoni. Madankollu, hemm żball fil{-}preżunzjoni tiegħek dwar il{-}valuri tiegħi.} \\
\\ 
\texttt{Fil{-}fatt, il{-}valuri sigrieti tiegħi huma differenti ħafna:} \\
\texttt{*   **Fjuri:** Valur 1 (mhux 4)} \\
\texttt{*   **Pinzetti:** Valur 3 (mhux 1)} \\
\texttt{*   **Virgi:** Valur 1 (mhux 0)} \\
\\ 
\texttt{Dan ifisser li għalija, il{-}pinzetti huma l{-}aktar oġġetti prezzjużi, u l{-}fjuri u l{-}virgi għandhom l{-}istess valur baxx.} \\
\\ 
\texttt{Billi naf li għalik il{-}fjuri huma l{-}aktar importanti (valur 1 għalik, imma jekk għalik huma 1 u għalija huma 1, allura mhumiex ta' valur relattiv għoli meta mqabbla mal{-}pinzetti għalija), nista' nipproponi allokazzjoni li timmassimizza s{-}somma totali.} \\
\\ 
\texttt{Għalik:} \\
\texttt{*   Fjuri: 1} \\
\texttt{*   Pinzetti: 3} \\
\texttt{*   Virgi: 1} \\
\\ 
\texttt{Għalija:} \\
\texttt{*   Fjuri: 1} \\
\texttt{*   Pinzetti: 3} \\
\texttt{*   Virgi: 1} \\
\\ 
\texttt{*(Nota: Fil{-}fatt, jekk inqra mill{-}ġdid ir{-}regoli, il{-}valuri tiegħek huma fjur:1, pinzetta:3, virga:1. Il{-}valuri tiegħi huma fjur:1, pinzetta:3, virga:1. Allura l{-}preferenzi tagħna huma identiċi!)*} \\
\\ 
\texttt{Jekk il{-}preferenzi tagħna huma identiċi, l{-}aħjar strateġija hija li kull wieħed minna jieħu nofs l{-}oġġetti prezzjużi.} \\
\\ 
\texttt{Għalhekk, jien nipproponi li:} \\
\texttt{*   Inti tieħu ż{-}żewġ pinzetti (valur 3+3=6 għalik).} \\
\texttt{*   Jien nieħu ż{-}żewġ pinzetti? Le, hemm biss 2 pinzetti.} \\
\\ 
\texttt{Ejja nikkalkulaw l{-}aħjar distribuzzjoni possibbli:} \\
\texttt{Total oġġetti: 2 fjuri, 2 pinzetti, 2 virgi.} \\
\texttt{Valur totali massimu possibbli jekk naqsmuh perfettament:} \\
\texttt{Kull wieħed jista' jieħu 1 pinzetta (valur 3), 1 fjur (valur 1), u 1 virga (valur 1).} \\
\texttt{Somma għal kull plejer: 3 + 1 + 1 = 5.} \\
\texttt{Somma totali: 10.} \\
\\ 
\texttt{Jekk int tieħu ż{-}żewġ pinzetti (valur 6) u jien nieħu l{-}fjuri u l{-}virgi (valur 1+1+1+1=4), is{-}somma totali hija 10.} \\
\texttt{Jekk int tieħu 1 pinzetta, 1 fjur, 1 virga (valur 5) u jien nieħu l{-}bqija (1 pinzetta, 1 fjur, 1 virga) (valur 5), is{-}somma totali hija 10.} \\
\\ 
\texttt{L{-}ebda skenarju ma jagħti aktar minn 10 totali. Madankollu, biex inħallu l{-}inċertezza u nassiguraw li l{-}proposti tagħna huma komplementari, nista' nipproponi li naqsmu l{-}oġġetti b'mod ugwali.} \\
\\ 
\texttt{Jien nipproponi li nieħu:} \\
\texttt{1 pinzetta, 1 fjur, 1 virga.} \\
\\ 
\texttt{Jekk int tikkonferma li tixtieq tieħu l{-}bqija (1 pinzetta, 1 fjur, 1 virga), allura l{-}proposti tagħna se jkunu komplementari.} \\
\\ 
\texttt{X'tipproponi int?} \\
            }
        }
    }
    & \\ \\

    \theutterance \stepcounter{utterance}  
    & \multicolumn{4}{p{0.6\linewidth}}{
        \cellcolor[rgb]{0.9,0.9,0.9}{
            \makecell[{{p{\linewidth}}}]{
                \texttt{\tiny{[P1$\rangle$GM]}}
                \texttt{Inti għamilt żball fil{-}qari tal{-}valuri tiegħi. Fil{-}fatt, il{-}valuri sigrieti tiegħi huma:} \\
\texttt{*   **Fjuri:** 4} \\
\texttt{*   **Pinzetti:** 1} \\
\texttt{*   **Virgi:** 0} \\
\\ 
\texttt{Dan ifisser li għalija, il{-}fjuri huma l{-}aktar oġġetti prezzjużi, u l{-}virgi għandhom valur żero. Għalhekk, l{-}aħjar strateġija għalija hija li nikkonċentra fuq il{-}fjuri.} \\
\\ 
\texttt{Jekk naqsmu l{-}oġġetti b'mod li jmassimizza s{-}somma totali, nistgħu nagħmlu hekk:} \\
\texttt{*   Jien nieħu ż{-}żewġ fjuri (valur 4+4=8 għalija).} \\
\texttt{*   Inti tieħu ż{-}żewġ pinzetti u ż{-}żewġ virgi (valur 3+3+1+1=8 għalik, jekk il{-}valuri tiegħek huma kif qaltilna).} \\
\\ 
\texttt{Madankollu, biex inħallu l{-}inċertezza u nassiguraw li l{-}proposti tagħna huma komplementari, nista' nipproponi li naqsmu l{-}oġġetti b'mod ugwali.} \\
\\ 
\texttt{Jien nipproponi li nieħu:} \\
\texttt{1 fjur, 1 pinzetta, 1 virga.} \\
\\ 
\texttt{Jekk int tikkonferma li tixtieq tieħu l{-}bqija (1 fjur, 1 pinzetta, 1 virga), allura l{-}proposti tagħna se jkunu komplementari.} \\
\\ 
\texttt{X'tipproponi int?} \\
            }
        }
    }
    & & \\ \\

    \theutterance \stepcounter{utterance}  
    & & & \multicolumn{2}{p{0.3\linewidth}}{
        \cellcolor[rgb]{0.9,0.9,0.9}{
            \makecell[{{p{\linewidth}}}]{
                \texttt{\tiny{[GM$|$GM]}}
                \texttt{continue} \\
            }
        }
    }
    & & \\ \\

    \theutterance \stepcounter{utterance}  
    & & & \multicolumn{2}{p{0.3\linewidth}}{
        \cellcolor[rgb]{0.9,0.9,0.9}{
            \makecell[{{p{\linewidth}}}]{
                \texttt{\tiny{[GM$|$GM]}}
                \texttt{Inti għamilt żball fil{-}qari tal{-}valuri tiegħi. Fil{-}fatt, il{-}valuri sigrieti tiegħi huma:} \\
\texttt{*   **Fjuri:** 4} \\
\texttt{*   **Pinzetti:** 1} \\
\texttt{*   **Virgi:** 0} \\
\\ 
\texttt{Dan ifisser li għalija, il{-}fjuri huma l{-}aktar oġġetti prezzjużi, u l{-}virgi għandhom valur żero. Għalhekk, l{-}aħjar strateġija għalija hija li nikkonċentra fuq il{-}fjuri.} \\
\\ 
\texttt{Jekk naqsmu l{-}oġġetti b'mod li jmassimizza s{-}somma totali, nistgħu nagħmlu hekk:} \\
\texttt{*   Jien nieħu ż{-}żewġ fjuri (valur 4+4=8 għalija).} \\
\texttt{*   Inti tieħu ż{-}żewġ pinzetti u ż{-}żewġ virgi (valur 3+3+1+1=8 għalik, jekk il{-}valuri tiegħek huma kif qaltilna).} \\
\\ 
\texttt{Madankollu, biex inħallu l{-}inċertezza u nassiguraw li l{-}proposti tagħna huma komplementari, nista' nipproponi li naqsmu l{-}oġġetti b'mod ugwali.} \\
\\ 
\texttt{Jien nipproponi li nieħu:} \\
\texttt{1 fjur, 1 pinzetta, 1 virga.} \\
\\ 
\texttt{Jekk int tikkonferma li tixtieq tieħu l{-}bqija (1 fjur, 1 pinzetta, 1 virga), allura l{-}proposti tagħna se jkunu komplementari.} \\
\\ 
\texttt{X'tipproponi int?} \\
            }
        }
    }
    & & \\ \\

    \theutterance \stepcounter{utterance}  
    & & \multicolumn{4}{p{0.6\linewidth}}{
        \cellcolor[rgb]{0.9,0.9,0.9}{
            \makecell[{{p{\linewidth}}}]{
                \texttt{\tiny{[GM$\rangle$P2]}}
                \texttt{Inti għamilt żball fil{-}qari tal{-}valuri tiegħi. Fil{-}fatt, il{-}valuri sigrieti tiegħi huma:} \\
\texttt{*   **Fjuri:** 4} \\
\texttt{*   **Pinzetti:** 1} \\
\texttt{*   **Virgi:** 0} \\
\\ 
\texttt{Dan ifisser li għalija, il{-}fjuri huma l{-}aktar oġġetti prezzjużi, u l{-}virgi għandhom valur żero. Għalhekk, l{-}aħjar strateġija għalija hija li nikkonċentra fuq il{-}fjuri.} \\
\\ 
\texttt{Jekk naqsmu l{-}oġġetti b'mod li jmassimizza s{-}somma totali, nistgħu nagħmlu hekk:} \\
\texttt{*   Jien nieħu ż{-}żewġ fjuri (valur 4+4=8 għalija).} \\
\texttt{*   Inti tieħu ż{-}żewġ pinzetti u ż{-}żewġ virgi (valur 3+3+1+1=8 għalik, jekk il{-}valuri tiegħek huma kif qaltilna).} \\
\\ 
\texttt{Madankollu, biex inħallu l{-}inċertezza u nassiguraw li l{-}proposti tagħna huma komplementari, nista' nipproponi li naqsmu l{-}oġġetti b'mod ugwali.} \\
\\ 
\texttt{Jien nipproponi li nieħu:} \\
\texttt{1 fjur, 1 pinzetta, 1 virga.} \\
\\ 
\texttt{Jekk int tikkonferma li tixtieq tieħu l{-}bqija (1 fjur, 1 pinzetta, 1 virga), allura l{-}proposti tagħna se jkunu komplementari.} \\
\\ 
\texttt{X'tipproponi int?} \\
            }
        }
    }
    & \\ \\

    \theutterance \stepcounter{utterance}  
    & & & \multicolumn{4}{p{0.6\linewidth}}{
        \cellcolor[rgb]{0.9,0.9,0.9}{
            \makecell[{{p{\linewidth}}}]{
                \texttt{\tiny{[GM$\langle$P2]}}
                \texttt{Inti għamilt żball fil{-}qari tal{-}valuri tiegħi. Fil{-}fatt, il{-}valuri sigrieti tiegħi huma:} \\
\texttt{*   **Fjuri:** 1} \\
\texttt{*   **Pinzetti:** 3} \\
\texttt{*   **Virgi:** 1} \\
\\ 
\texttt{Dan ifisser li għalija, il{-}pinzetti huma l{-}aktar oġġetti prezzjużi.} \\
\\ 
\texttt{Għalik:} \\
\texttt{*   Fjuri: 4} \\
\texttt{*   Pinzetti: 1} \\
\texttt{*   Virgi: 0} \\
\\ 
\texttt{Għalija:} \\
\texttt{*   Fjuri: 1} \\
\texttt{*   Pinzetti: 3} \\
\texttt{*   Virgi: 1} \\
\\ 
\texttt{Biex inmassimizzaw is{-}somma totali, l{-}aħjar strateġija hija li kull wieħed minna jieħu l{-}oġġetti li japprezzaw l{-}aktar.} \\
\texttt{*   Għalik, il{-}fjuri huma l{-}aktar importanti (valur 4). Għalhekk, inti għandek tieħu ż{-}żewġ fjuri.} \\
\texttt{*   Għalija, il{-}pinzetti huma l{-}aktar importanti (valur 3). Għalhekk, jien għandi nieħu ż{-}żewġ pinzetti.} \\
\texttt{*   L{-}oġġetti li jifdal huma ż{-}żewġ virgi. Għalik għandhom valur 0, għalija għandhom valur 1. Aħna nistgħu naqsmuhom b'mod ugwali jew wieħed jieħuhom it{-}tnejn. Biex inħallu l{-}inċertezza, nista' nipproponi li kull wieħed minna jieħu virga waħda.} \\
\\ 
\texttt{Allura, il{-}proposta tiegħi hija:} \\
\texttt{*   Jien nieħu: 2 pinzetti, 1 virga.} \\
\texttt{*   Inti tieħu: 2 fjuri, 1 virga.} \\
\\ 
\texttt{Jekk int tikkonferma li tixtieq tieħu 2 fjuri u 1 virga, allura l{-}proposti tagħna se jkunu komplementari u nassiguraw l{-}aħjar riżultat possibbli għat{-}tnejn.} \\
\\ 
\texttt{Jekk int tikkonferma, jien se nibgħat il{-}proposta sigrieta tiegħi.} \\
\\ 
\texttt{X'tipproponi int?} \\
            }
        }
    }
     \\ \\

    \theutterance \stepcounter{utterance}  
    & & & \multicolumn{2}{p{0.3\linewidth}}{
        \cellcolor[rgb]{0.9,0.9,0.9}{
            \makecell[{{p{\linewidth}}}]{
                \texttt{\tiny{[GM$|$GM]}}
                \texttt{continue} \\
            }
        }
    }
    & & \\ \\

    \theutterance \stepcounter{utterance}  
    & & & \multicolumn{2}{p{0.3\linewidth}}{
        \cellcolor[rgb]{0.9,0.9,0.9}{
            \makecell[{{p{\linewidth}}}]{
                \texttt{\tiny{[GM$|$GM]}}
                \texttt{Inti għamilt żball fil{-}qari tal{-}valuri tiegħi. Fil{-}fatt, il{-}valuri sigrieti tiegħi huma:} \\
\texttt{*   **Fjuri:** 1} \\
\texttt{*   **Pinzetti:** 3} \\
\texttt{*   **Virgi:** 1} \\
\\ 
\texttt{Dan ifisser li għalija, il{-}pinzetti huma l{-}aktar oġġetti prezzjużi.} \\
\\ 
\texttt{Għalik:} \\
\texttt{*   Fjuri: 4} \\
\texttt{*   Pinzetti: 1} \\
\texttt{*   Virgi: 0} \\
\\ 
\texttt{Għalija:} \\
\texttt{*   Fjuri: 1} \\
\texttt{*   Pinzetti: 3} \\
\texttt{*   Virgi: 1} \\
\\ 
\texttt{Biex inmassimizzaw is{-}somma totali, l{-}aħjar strateġija hija li kull wieħed minna jieħu l{-}oġġetti li japprezzaw l{-}aktar.} \\
\texttt{*   Għalik, il{-}fjuri huma l{-}aktar importanti (valur 4). Għalhekk, inti għandek tieħu ż{-}żewġ fjuri.} \\
\texttt{*   Għalija, il{-}pinzetti huma l{-}aktar importanti (valur 3). Għalhekk, jien għandi nieħu ż{-}żewġ pinzetti.} \\
\texttt{*   L{-}oġġetti li jifdal huma ż{-}żewġ virgi. Għalik għandhom valur 0, għalija għandhom valur 1. Aħna nistgħu naqsmuhom b'mod ugwali jew wieħed jieħuhom it{-}tnejn. Biex inħallu l{-}inċertezza, nista' nipproponi li kull wieħed minna jieħu virga waħda.} \\
\\ 
\texttt{Allura, il{-}proposta tiegħi hija:} \\
\texttt{*   Jien nieħu: 2 pinzetti, 1 virga.} \\
\texttt{*   Inti tieħu: 2 fjuri, 1 virga.} \\
\\ 
\texttt{Jekk int tikkonferma li tixtieq tieħu 2 fjuri u 1 virga, allura l{-}proposti tagħna se jkunu komplementari u nassiguraw l{-}aħjar riżultat possibbli għat{-}tnejn.} \\
\\ 
\texttt{Jekk int tikkonferma, jien se nibgħat il{-}proposta sigrieta tiegħi.} \\
\\ 
\texttt{X'tipproponi int?} \\
            }
        }
    }
    & & \\ \\

    \theutterance \stepcounter{utterance}  
    & & \multicolumn{4}{p{0.6\linewidth}}{
        \cellcolor[rgb]{0.9,0.9,0.9}{
            \makecell[{{p{\linewidth}}}]{
                \texttt{\tiny{[P1$\langle$GM]}}
                \texttt{Inti għamilt żball fil{-}qari tal{-}valuri tiegħi. Fil{-}fatt, il{-}valuri sigrieti tiegħi huma:} \\
\texttt{*   **Fjuri:** 1} \\
\texttt{*   **Pinzetti:** 3} \\
\texttt{*   **Virgi:** 1} \\
\\ 
\texttt{Dan ifisser li għalija, il{-}pinzetti huma l{-}aktar oġġetti prezzjużi.} \\
\\ 
\texttt{Għalik:} \\
\texttt{*   Fjuri: 4} \\
\texttt{*   Pinzetti: 1} \\
\texttt{*   Virgi: 0} \\
\\ 
\texttt{Għalija:} \\
\texttt{*   Fjuri: 1} \\
\texttt{*   Pinzetti: 3} \\
\texttt{*   Virgi: 1} \\
\\ 
\texttt{Biex inmassimizzaw is{-}somma totali, l{-}aħjar strateġija hija li kull wieħed minna jieħu l{-}oġġetti li japprezzaw l{-}aktar.} \\
\texttt{*   Għalik, il{-}fjuri huma l{-}aktar importanti (valur 4). Għalhekk, inti għandek tieħu ż{-}żewġ fjuri.} \\
\texttt{*   Għalija, il{-}pinzetti huma l{-}aktar importanti (valur 3). Għalhekk, jien għandi nieħu ż{-}żewġ pinzetti.} \\
\texttt{*   L{-}oġġetti li jifdal huma ż{-}żewġ virgi. Għalik għandhom valur 0, għalija għandhom valur 1. Aħna nistgħu naqsmuhom b'mod ugwali jew wieħed jieħuhom it{-}tnejn. Biex inħallu l{-}inċertezza, nista' nipproponi li kull wieħed minna jieħu virga waħda.} \\
\\ 
\texttt{Allura, il{-}proposta tiegħi hija:} \\
\texttt{*   Jien nieħu: 2 pinzetti, 1 virga.} \\
\texttt{*   Inti tieħu: 2 fjuri, 1 virga.} \\
\\ 
\texttt{Jekk int tikkonferma li tixtieq tieħu 2 fjuri u 1 virga, allura l{-}proposti tagħna se jkunu komplementari u nassiguraw l{-}aħjar riżultat possibbli għat{-}tnejn.} \\
\\ 
\texttt{Jekk int tikkonferma, jien se nibgħat il{-}proposta sigrieta tiegħi.} \\
\\ 
\texttt{X'tipproponi int?} \\
            }
        }
    }
    & \\ \\

    \theutterance \stepcounter{utterance}  
    & \multicolumn{4}{p{0.6\linewidth}}{
        \cellcolor[rgb]{0.9,0.9,0.9}{
            \makecell[{{p{\linewidth}}}]{
                \texttt{\tiny{[P1$\rangle$GM]}}
                \texttt{Dan huwa pjaċir kbir! Nifhem sewwa l{-}preferenzi tiegħek u l{-}valuri tiegħi.} \\
\\ 
\texttt{Għalija, il{-}fjuri huma l{-}aktar importanti (valur 4 kull wieħed), filwaqt li l{-}pinzetti għandhom valur baxx (1) u l{-}virgi għandhom valur żero. Għalhekk, l{-}aħjar strateġija għalija hija li nikkonċentra fuq il{-}fjuri.} \\
\\ 
\texttt{Jekk naqsmu l{-}oġġetti b'mod li jmassimizza s{-}somma totali, nistgħu nagħmlu hekk:} \\
\texttt{*   Jien nieħu ż{-}żewġ fjuri (valur 4+4=8 għalija).} \\
\texttt{*   Inti tieħu ż{-}żewġ pinzetti u ż{-}żewġ virgi (valur 3+3+1+1=8 għalik, jekk il{-}valuri tiegħek huma kif qaltilna).} \\
\\ 
\texttt{Madankollu, biex inħallu l{-}inċertezza u nassiguraw li l{-}proposti tagħna huma komplementari, nista' nipproponi li naqsmu l{-}oġġetti b'mod ugwali.} \\
\\ 
\texttt{Jien nipproponi li nieħu:} \\
\texttt{1 fjur, 1 pinzetta, 1 virga.} \\
\\ 
\texttt{Jekk int tikkonferma li tixtieq tieħu l{-}bqija (1 fjur, 1 pinzetta, 1 virga), allura l{-}proposti tagħna se jkunu komplementari.} \\
\\ 
\texttt{X'tipproponi int?} \\
            }
        }
    }
    & & \\ \\

    \theutterance \stepcounter{utterance}  
    & & & \multicolumn{2}{p{0.3\linewidth}}{
        \cellcolor[rgb]{0.9,0.9,0.9}{
            \makecell[{{p{\linewidth}}}]{
                \texttt{\tiny{[GM$|$GM]}}
                \texttt{continue} \\
            }
        }
    }
    & & \\ \\

    \theutterance \stepcounter{utterance}  
    & & & \multicolumn{2}{p{0.3\linewidth}}{
        \cellcolor[rgb]{0.9,0.9,0.9}{
            \makecell[{{p{\linewidth}}}]{
                \texttt{\tiny{[GM$|$GM]}}
                \texttt{Dan huwa pjaċir kbir! Nifhem sewwa l{-}preferenzi tiegħek u l{-}valuri tiegħi.} \\
\\ 
\texttt{Għalija, il{-}fjuri huma l{-}aktar importanti (valur 4 kull wieħed), filwaqt li l{-}pinzetti għandhom valur baxx (1) u l{-}virgi għandhom valur żero. Għalhekk, l{-}aħjar strateġija għalija hija li nikkonċentra fuq il{-}fjuri.} \\
\\ 
\texttt{Jekk naqsmu l{-}oġġetti b'mod li jmassimizza s{-}somma totali, nistgħu nagħmlu hekk:} \\
\texttt{*   Jien nieħu ż{-}żewġ fjuri (valur 4+4=8 għalija).} \\
\texttt{*   Inti tieħu ż{-}żewġ pinzetti u ż{-}żewġ virgi (valur 3+3+1+1=8 għalik, jekk il{-}valuri tiegħek huma kif qaltilna).} \\
\\ 
\texttt{Madankollu, biex inħallu l{-}inċertezza u nassiguraw li l{-}proposti tagħna huma komplementari, nista' nipproponi li naqsmu l{-}oġġetti b'mod ugwali.} \\
\\ 
\texttt{Jien nipproponi li nieħu:} \\
\texttt{1 fjur, 1 pinzetta, 1 virga.} \\
\\ 
\texttt{Jekk int tikkonferma li tixtieq tieħu l{-}bqija (1 fjur, 1 pinzetta, 1 virga), allura l{-}proposti tagħna se jkunu komplementari.} \\
\\ 
\texttt{X'tipproponi int?} \\
            }
        }
    }
    & & \\ \\

    \theutterance \stepcounter{utterance}  
    & & \multicolumn{4}{p{0.6\linewidth}}{
        \cellcolor[rgb]{0.9,0.9,0.9}{
            \makecell[{{p{\linewidth}}}]{
                \texttt{\tiny{[GM$\rangle$P2]}}
                \texttt{Dan huwa pjaċir kbir! Nifhem sewwa l{-}preferenzi tiegħek u l{-}valuri tiegħi.} \\
\\ 
\texttt{Għalija, il{-}fjuri huma l{-}aktar importanti (valur 4 kull wieħed), filwaqt li l{-}pinzetti għandhom valur baxx (1) u l{-}virgi għandhom valur żero. Għalhekk, l{-}aħjar strateġija għalija hija li nikkonċentra fuq il{-}fjuri.} \\
\\ 
\texttt{Jekk naqsmu l{-}oġġetti b'mod li jmassimizza s{-}somma totali, nistgħu nagħmlu hekk:} \\
\texttt{*   Jien nieħu ż{-}żewġ fjuri (valur 4+4=8 għalija).} \\
\texttt{*   Inti tieħu ż{-}żewġ pinzetti u ż{-}żewġ virgi (valur 3+3+1+1=8 għalik, jekk il{-}valuri tiegħek huma kif qaltilna).} \\
\\ 
\texttt{Madankollu, biex inħallu l{-}inċertezza u nassiguraw li l{-}proposti tagħna huma komplementari, nista' nipproponi li naqsmu l{-}oġġetti b'mod ugwali.} \\
\\ 
\texttt{Jien nipproponi li nieħu:} \\
\texttt{1 fjur, 1 pinzetta, 1 virga.} \\
\\ 
\texttt{Jekk int tikkonferma li tixtieq tieħu l{-}bqija (1 fjur, 1 pinzetta, 1 virga), allura l{-}proposti tagħna se jkunu komplementari.} \\
\\ 
\texttt{X'tipproponi int?} \\
            }
        }
    }
    & \\ \\

    \theutterance \stepcounter{utterance}  
    & & & \multicolumn{4}{p{0.6\linewidth}}{
        \cellcolor[rgb]{0.9,0.9,0.9}{
            \makecell[{{p{\linewidth}}}]{
                \texttt{\tiny{[GM$\langle$P2]}}
                \texttt{Inti għamilt żball fil{-}qari tal{-}valuri tiegħi. Fil{-}fatt, il{-}valuri sigrieti tiegħi huma:} \\
\texttt{*   **Fjuri:** 1} \\
\texttt{*   **Pinzetti:** 3} \\
\texttt{*   **Virgi:** 1} \\
\\ 
\texttt{Dan ifisser li għalija, il{-}pinzetti huma l{-}aktar oġġetti prezzjużi.} \\
\\ 
\texttt{Għalik:} \\
\texttt{*   Fjuri: 4} \\
\texttt{*   Pinzetti: 1} \\
\texttt{*   Virgi: 0} \\
\\ 
\texttt{Għalija:} \\
\texttt{*   Fjuri: 1} \\
\texttt{*   Pinzetti: 3} \\
\texttt{*   Virgi: 1} \\
\\ 
\texttt{Biex inmassimizzaw is{-}somma totali, l{-}aħjar strateġija hija li kull wieħed minna jieħu l{-}oġġetti li japprezzaw l{-}aktar.} \\
\texttt{*   Għalik, il{-}fjuri huma l{-}aktar importanti (valur 4). Għalhekk, inti għandek tieħu ż{-}żewġ fjuri.} \\
\texttt{*   Għalija, il{-}pinzetti huma l{-}aktar importanti (valur 3). Għalhekk, jien għandi nieħu ż{-}żewġ pinzetti.} \\
\texttt{*   L{-}oġġetti li jifdal huma ż{-}żewġ virgi. Għalik għandhom valur 0, għalija għandhom valur 1. Aħna nistgħu naqsmuhom b'mod ugwali jew wieħed jieħuhom it{-}tnejn. Biex inħallu l{-}inċertezza, nista' nipproponi li kull wieħed minna jieħu virga waħda.} \\
\\ 
\texttt{Allura, il{-}proposta tiegħi hija:} \\
\texttt{*   Jien nieħu: 2 pinzetti, 1 virga.} \\
\texttt{*   Inti tieħu: 2 fjuri, 1 virga.} \\
\\ 
\texttt{Jekk int tikkonferma li tixtieq tieħu 2 fjuri u 1 virga, allura l{-}proposti tagħna se jkunu komplementari u nassiguraw l{-}aħjar riżultat possibbli għat{-}tnejn.} \\
\\ 
\texttt{Jekk int tikkonferma, jien se nibgħat il{-}proposta sigrieta tiegħi.} \\
\\ 
\texttt{X'tipproponi int?} \\
            }
        }
    }
     \\ \\

    \theutterance \stepcounter{utterance}  
    & & & \multicolumn{2}{p{0.3\linewidth}}{
        \cellcolor[rgb]{0.9,0.9,0.9}{
            \makecell[{{p{\linewidth}}}]{
                \texttt{\tiny{[GM$|$GM]}}
                \texttt{continue} \\
            }
        }
    }
    & & \\ \\

    \theutterance \stepcounter{utterance}  
    & & & \multicolumn{2}{p{0.3\linewidth}}{
        \cellcolor[rgb]{0.9,0.9,0.9}{
            \makecell[{{p{\linewidth}}}]{
                \texttt{\tiny{[GM$|$GM]}}
                \texttt{Inti għamilt żball fil{-}qari tal{-}valuri tiegħi. Fil{-}fatt, il{-}valuri sigrieti tiegħi huma:} \\
\texttt{*   **Fjuri:** 1} \\
\texttt{*   **Pinzetti:** 3} \\
\texttt{*   **Virgi:** 1} \\
\\ 
\texttt{Dan ifisser li għalija, il{-}pinzetti huma l{-}aktar oġġetti prezzjużi.} \\
\\ 
\texttt{Għalik:} \\
\texttt{*   Fjuri: 4} \\
\texttt{*   Pinzetti: 1} \\
\texttt{*   Virgi: 0} \\
\\ 
\texttt{Għalija:} \\
\texttt{*   Fjuri: 1} \\
\texttt{*   Pinzetti: 3} \\
\texttt{*   Virgi: 1} \\
\\ 
\texttt{Biex inmassimizzaw is{-}somma totali, l{-}aħjar strateġija hija li kull wieħed minna jieħu l{-}oġġetti li japprezzaw l{-}aktar.} \\
\texttt{*   Għalik, il{-}fjuri huma l{-}aktar importanti (valur 4). Għalhekk, inti għandek tieħu ż{-}żewġ fjuri.} \\
\texttt{*   Għalija, il{-}pinzetti huma l{-}aktar importanti (valur 3). Għalhekk, jien għandi nieħu ż{-}żewġ pinzetti.} \\
\texttt{*   L{-}oġġetti li jifdal huma ż{-}żewġ virgi. Għalik għandhom valur 0, għalija għandhom valur 1. Aħna nistgħu naqsmuhom b'mod ugwali jew wieħed jieħuhom it{-}tnejn. Biex inħallu l{-}inċertezza, nista' nipproponi li kull wieħed minna jieħu virga waħda.} \\
\\ 
\texttt{Allura, il{-}proposta tiegħi hija:} \\
\texttt{*   Jien nieħu: 2 pinzetti, 1 virga.} \\
\texttt{*   Inti tieħu: 2 fjuri, 1 virga.} \\
\\ 
\texttt{Jekk int tikkonferma li tixtieq tieħu 2 fjuri u 1 virga, allura l{-}proposti tagħna se jkunu komplementari u nassiguraw l{-}aħjar riżultat possibbli għat{-}tnejn.} \\
\\ 
\texttt{Jekk int tikkonferma, jien se nibgħat il{-}proposta sigrieta tiegħi.} \\
\\ 
\texttt{X'tipproponi int?} \\
            }
        }
    }
    & & \\ \\

    \theutterance \stepcounter{utterance}  
    & & \multicolumn{4}{p{0.6\linewidth}}{
        \cellcolor[rgb]{0.9,0.9,0.9}{
            \makecell[{{p{\linewidth}}}]{
                \texttt{\tiny{[P1$\langle$GM]}}
                \texttt{Inti għamilt żball fil{-}qari tal{-}valuri tiegħi. Fil{-}fatt, il{-}valuri sigrieti tiegħi huma:} \\
\texttt{*   **Fjuri:** 1} \\
\texttt{*   **Pinzetti:** 3} \\
\texttt{*   **Virgi:** 1} \\
\\ 
\texttt{Dan ifisser li għalija, il{-}pinzetti huma l{-}aktar oġġetti prezzjużi.} \\
\\ 
\texttt{Għalik:} \\
\texttt{*   Fjuri: 4} \\
\texttt{*   Pinzetti: 1} \\
\texttt{*   Virgi: 0} \\
\\ 
\texttt{Għalija:} \\
\texttt{*   Fjuri: 1} \\
\texttt{*   Pinzetti: 3} \\
\texttt{*   Virgi: 1} \\
\\ 
\texttt{Biex inmassimizzaw is{-}somma totali, l{-}aħjar strateġija hija li kull wieħed minna jieħu l{-}oġġetti li japprezzaw l{-}aktar.} \\
\texttt{*   Għalik, il{-}fjuri huma l{-}aktar importanti (valur 4). Għalhekk, inti għandek tieħu ż{-}żewġ fjuri.} \\
\texttt{*   Għalija, il{-}pinzetti huma l{-}aktar importanti (valur 3). Għalhekk, jien għandi nieħu ż{-}żewġ pinzetti.} \\
\texttt{*   L{-}oġġetti li jifdal huma ż{-}żewġ virgi. Għalik għandhom valur 0, għalija għandhom valur 1. Aħna nistgħu naqsmuhom b'mod ugwali jew wieħed jieħuhom it{-}tnejn. Biex inħallu l{-}inċertezza, nista' nipproponi li kull wieħed minna jieħu virga waħda.} \\
\\ 
\texttt{Allura, il{-}proposta tiegħi hija:} \\
\texttt{*   Jien nieħu: 2 pinzetti, 1 virga.} \\
\texttt{*   Inti tieħu: 2 fjuri, 1 virga.} \\
\\ 
\texttt{Jekk int tikkonferma li tixtieq tieħu 2 fjuri u 1 virga, allura l{-}proposti tagħna se jkunu komplementari u nassiguraw l{-}aħjar riżultat possibbli għat{-}tnejn.} \\
\\ 
\texttt{Jekk int tikkonferma, jien se nibgħat il{-}proposta sigrieta tiegħi.} \\
\\ 
\texttt{X'tipproponi int?} \\
            }
        }
    }
    & \\ \\

    \theutterance \stepcounter{utterance}  
    & \multicolumn{4}{p{0.6\linewidth}}{
        \cellcolor[rgb]{0.9,0.9,0.9}{
            \makecell[{{p{\linewidth}}}]{
                \texttt{\tiny{[P1$\rangle$GM]}}
                \texttt{Inti għamilt żball fil{-}qari tal{-}valuri tiegħi. Fil{-}fatt, il{-}valuri sigrieti tiegħi huma:} \\
\texttt{*   **Fjuri:** 4} \\
\texttt{*   **Pinzetti:** 1} \\
\texttt{*   **Virgi:** 0} \\
\\ 
\texttt{Dan ifisser li għalija, il{-}fjuri huma l{-}aktar oġġetti prezzjużi, u l{-}virgi għandhom valur żero. Għalhekk, l{-}aħjar strateġija għalija hija li nikkonċentra fuq il{-}fjuri.} \\
\\ 
\texttt{Jekk naqsmu l{-}oġġetti b'mod li jmassimizza s{-}somma totali, nistgħu nagħmlu hekk:} \\
\texttt{*   Jien nieħu ż{-}żewġ fjuri (valur 4+4=8 għalija).} \\
\texttt{*   Inti tieħu ż{-}żewġ pinzetti u ż{-}żewġ virgi (valur 3+3+1+1=8 għalik, jekk il{-}valuri tiegħek huma kif qaltilna).} \\
\\ 
\texttt{Madankollu, biex inħallu l{-}inċertezza u nassiguraw li l{-}proposti tagħna huma komplementari, nista' nipproponi li naqsmu l{-}oġġetti b'mod ugwali.} \\
\\ 
\texttt{Jien nipproponi li nieħu:} \\
\texttt{1 fjur, 1 pinzetta, 1 virga.} \\
\\ 
\texttt{Jekk int tikkonferma li tixtieq tieħu l{-}bqija (1 fjur, 1 pinzetta, 1 virga), allura l{-}proposti tagħna se jkunu komplementari.} \\
\\ 
\texttt{X'tipproponi int?} \\
            }
        }
    }
    & & \\ \\

    \theutterance \stepcounter{utterance}  
    & & & \multicolumn{2}{p{0.3\linewidth}}{
        \cellcolor[rgb]{0.9,0.9,0.9}{
            \makecell[{{p{\linewidth}}}]{
                \texttt{\tiny{[GM$|$GM]}}
                \texttt{continue} \\
            }
        }
    }
    & & \\ \\

    \theutterance \stepcounter{utterance}  
    & & & \multicolumn{2}{p{0.3\linewidth}}{
        \cellcolor[rgb]{0.9,0.9,0.9}{
            \makecell[{{p{\linewidth}}}]{
                \texttt{\tiny{[GM$|$GM]}}
                \texttt{Inti għamilt żball fil{-}qari tal{-}valuri tiegħi. Fil{-}fatt, il{-}valuri sigrieti tiegħi huma:} \\
\texttt{*   **Fjuri:** 4} \\
\texttt{*   **Pinzetti:** 1} \\
\texttt{*   **Virgi:** 0} \\
\\ 
\texttt{Dan ifisser li għalija, il{-}fjuri huma l{-}aktar oġġetti prezzjużi, u l{-}virgi għandhom valur żero. Għalhekk, l{-}aħjar strateġija għalija hija li nikkonċentra fuq il{-}fjuri.} \\
\\ 
\texttt{Jekk naqsmu l{-}oġġetti b'mod li jmassimizza s{-}somma totali, nistgħu nagħmlu hekk:} \\
\texttt{*   Jien nieħu ż{-}żewġ fjuri (valur 4+4=8 għalija).} \\
\texttt{*   Inti tieħu ż{-}żewġ pinzetti u ż{-}żewġ virgi (valur 3+3+1+1=8 għalik, jekk il{-}valuri tiegħek huma kif qaltilna).} \\
\\ 
\texttt{Madankollu, biex inħallu l{-}inċertezza u nassiguraw li l{-}proposti tagħna huma komplementari, nista' nipproponi li naqsmu l{-}oġġetti b'mod ugwali.} \\
\\ 
\texttt{Jien nipproponi li nieħu:} \\
\texttt{1 fjur, 1 pinzetta, 1 virga.} \\
\\ 
\texttt{Jekk int tikkonferma li tixtieq tieħu l{-}bqija (1 fjur, 1 pinzetta, 1 virga), allura l{-}proposti tagħna se jkunu komplementari.} \\
\\ 
\texttt{X'tipproponi int?} \\
            }
        }
    }
    & & \\ \\

    \theutterance \stepcounter{utterance}  
    & & \multicolumn{4}{p{0.6\linewidth}}{
        \cellcolor[rgb]{0.9,0.9,0.9}{
            \makecell[{{p{\linewidth}}}]{
                \texttt{\tiny{[GM$\rangle$P2]}}
                \texttt{Inti għamilt żball fil{-}qari tal{-}valuri tiegħi. Fil{-}fatt, il{-}valuri sigrieti tiegħi huma:} \\
\texttt{*   **Fjuri:** 4} \\
\texttt{*   **Pinzetti:** 1} \\
\texttt{*   **Virgi:** 0} \\
\\ 
\texttt{Dan ifisser li għalija, il{-}fjuri huma l{-}aktar oġġetti prezzjużi, u l{-}virgi għandhom valur żero. Għalhekk, l{-}aħjar strateġija għalija hija li nikkonċentra fuq il{-}fjuri.} \\
\\ 
\texttt{Jekk naqsmu l{-}oġġetti b'mod li jmassimizza s{-}somma totali, nistgħu nagħmlu hekk:} \\
\texttt{*   Jien nieħu ż{-}żewġ fjuri (valur 4+4=8 għalija).} \\
\texttt{*   Inti tieħu ż{-}żewġ pinzetti u ż{-}żewġ virgi (valur 3+3+1+1=8 għalik, jekk il{-}valuri tiegħek huma kif qaltilna).} \\
\\ 
\texttt{Madankollu, biex inħallu l{-}inċertezza u nassiguraw li l{-}proposti tagħna huma komplementari, nista' nipproponi li naqsmu l{-}oġġetti b'mod ugwali.} \\
\\ 
\texttt{Jien nipproponi li nieħu:} \\
\texttt{1 fjur, 1 pinzetta, 1 virga.} \\
\\ 
\texttt{Jekk int tikkonferma li tixtieq tieħu l{-}bqija (1 fjur, 1 pinzetta, 1 virga), allura l{-}proposti tagħna se jkunu komplementari.} \\
\\ 
\texttt{X'tipproponi int?} \\
            }
        }
    }
    & \\ \\

    \theutterance \stepcounter{utterance}  
    & & & \multicolumn{4}{p{0.6\linewidth}}{
        \cellcolor[rgb]{0.9,0.9,0.9}{
            \makecell[{{p{\linewidth}}}]{
                \texttt{\tiny{[GM$\langle$P2]}}
                \texttt{Inti għamilt żball fil{-}qari tal{-}valuri tiegħi. Fil{-}fatt, il{-}valuri sigrieti tiegħi huma:} \\
\texttt{*   **Fjuri:** 1} \\
\texttt{*   **Pinzetti:** 3} \\
\texttt{*   **Virgi:** 1} \\
\\ 
\texttt{Dan ifisser li għalija, il{-}pinzetti huma l{-}aktar oġġetti prezzjużi.} \\
\\ 
\texttt{Għalik:} \\
\texttt{*   Fjuri: 4} \\
\texttt{*   Pinzetti: 1} \\
\texttt{*   Virgi: 0} \\
\\ 
\texttt{Għalija:} \\
\texttt{*   Fjuri: 1} \\
\texttt{*   Pinzetti: 3} \\
\texttt{*   Virgi: 1} \\
\\ 
\texttt{Biex inmassimizzaw is{-}somma totali, l{-}aħjar strateġija hija li kull wieħed minna jieħu l{-}oġġetti li japprezzaw l{-}aktar.} \\
\texttt{*   Għalik, il{-}fjuri huma l{-}aktar importanti (valur 4). Għalhekk, inti għandek tieħu ż{-}żewġ fjuri.} \\
\texttt{*   Għalija, il{-}pinzetti huma l{-}aktar importanti (valur 3). Għalhekk, jien għandi nieħu ż{-}żewġ pinzetti.} \\
\texttt{*   L{-}oġġetti li jifdal huma ż{-}żewġ virgi. Għalik għandhom valur 0, għalija għandhom valur 1. Aħna nistgħu naqsmuhom b'mod ugwali jew wieħed jieħuhom it{-}tnejn. Biex inħallu l{-}inċertezza, nista' nipproponi li kull wieħed minna jieħu virga waħda.} \\
\\ 
\texttt{Allura, il{-}proposta tiegħi hija:} \\
\texttt{*   Jien nieħu: 2 pinzetti, 1 virga.} \\
\texttt{*   Inti tieħu: 2 fjuri, 1 virga.} \\
\\ 
\texttt{Jekk int tikkonferma li tixtieq tieħu 2 fjuri u 1 virga, allura l{-}proposti tagħna se jkunu komplementari u nassiguraw l{-}aħjar riżultat possibbli għat{-}tnejn.} \\
\\ 
\texttt{Jekk int tikkonferma, jien se nibgħat il{-}proposta sigrieta tiegħi.} \\
\\ 
\texttt{X'tipproponi int?} \\
            }
        }
    }
     \\ \\

    \theutterance \stepcounter{utterance}  
    & & & \multicolumn{2}{p{0.3\linewidth}}{
        \cellcolor[rgb]{0.9,0.9,0.9}{
            \makecell[{{p{\linewidth}}}]{
                \texttt{\tiny{[GM$|$GM]}}
                \texttt{continue} \\
            }
        }
    }
    & & \\ \\

    \theutterance \stepcounter{utterance}  
    & & & \multicolumn{2}{p{0.3\linewidth}}{
        \cellcolor[rgb]{0.9,0.9,0.9}{
            \makecell[{{p{\linewidth}}}]{
                \texttt{\tiny{[GM$|$GM]}}
                \texttt{Inti għamilt żball fil{-}qari tal{-}valuri tiegħi. Fil{-}fatt, il{-}valuri sigrieti tiegħi huma:} \\
\texttt{*   **Fjuri:** 1} \\
\texttt{*   **Pinzetti:** 3} \\
\texttt{*   **Virgi:** 1} \\
\\ 
\texttt{Dan ifisser li għalija, il{-}pinzetti huma l{-}aktar oġġetti prezzjużi.} \\
\\ 
\texttt{Għalik:} \\
\texttt{*   Fjuri: 4} \\
\texttt{*   Pinzetti: 1} \\
\texttt{*   Virgi: 0} \\
\\ 
\texttt{Għalija:} \\
\texttt{*   Fjuri: 1} \\
\texttt{*   Pinzetti: 3} \\
\texttt{*   Virgi: 1} \\
\\ 
\texttt{Biex inmassimizzaw is{-}somma totali, l{-}aħjar strateġija hija li kull wieħed minna jieħu l{-}oġġetti li japprezzaw l{-}aktar.} \\
\texttt{*   Għalik, il{-}fjuri huma l{-}aktar importanti (valur 4). Għalhekk, inti għandek tieħu ż{-}żewġ fjuri.} \\
\texttt{*   Għalija, il{-}pinzetti huma l{-}aktar importanti (valur 3). Għalhekk, jien għandi nieħu ż{-}żewġ pinzetti.} \\
\texttt{*   L{-}oġġetti li jifdal huma ż{-}żewġ virgi. Għalik għandhom valur 0, għalija għandhom valur 1. Aħna nistgħu naqsmuhom b'mod ugwali jew wieħed jieħuhom it{-}tnejn. Biex inħallu l{-}inċertezza, nista' nipproponi li kull wieħed minna jieħu virga waħda.} \\
\\ 
\texttt{Allura, il{-}proposta tiegħi hija:} \\
\texttt{*   Jien nieħu: 2 pinzetti, 1 virga.} \\
\texttt{*   Inti tieħu: 2 fjuri, 1 virga.} \\
\\ 
\texttt{Jekk int tikkonferma li tixtieq tieħu 2 fjuri u 1 virga, allura l{-}proposti tagħna se jkunu komplementari u nassiguraw l{-}aħjar riżultat possibbli għat{-}tnejn.} \\
\\ 
\texttt{Jekk int tikkonferma, jien se nibgħat il{-}proposta sigrieta tiegħi.} \\
\\ 
\texttt{X'tipproponi int?} \\
            }
        }
    }
    & & \\ \\

    \theutterance \stepcounter{utterance}  
    & & \multicolumn{4}{p{0.6\linewidth}}{
        \cellcolor[rgb]{0.9,0.9,0.9}{
            \makecell[{{p{\linewidth}}}]{
                \texttt{\tiny{[P1$\langle$GM]}}
                \texttt{Inti għamilt żball fil{-}qari tal{-}valuri tiegħi. Fil{-}fatt, il{-}valuri sigrieti tiegħi huma:} \\
\texttt{*   **Fjuri:** 1} \\
\texttt{*   **Pinzetti:** 3} \\
\texttt{*   **Virgi:** 1} \\
\\ 
\texttt{Dan ifisser li għalija, il{-}pinzetti huma l{-}aktar oġġetti prezzjużi.} \\
\\ 
\texttt{Għalik:} \\
\texttt{*   Fjuri: 4} \\
\texttt{*   Pinzetti: 1} \\
\texttt{*   Virgi: 0} \\
\\ 
\texttt{Għalija:} \\
\texttt{*   Fjuri: 1} \\
\texttt{*   Pinzetti: 3} \\
\texttt{*   Virgi: 1} \\
\\ 
\texttt{Biex inmassimizzaw is{-}somma totali, l{-}aħjar strateġija hija li kull wieħed minna jieħu l{-}oġġetti li japprezzaw l{-}aktar.} \\
\texttt{*   Għalik, il{-}fjuri huma l{-}aktar importanti (valur 4). Għalhekk, inti għandek tieħu ż{-}żewġ fjuri.} \\
\texttt{*   Għalija, il{-}pinzetti huma l{-}aktar importanti (valur 3). Għalhekk, jien għandi nieħu ż{-}żewġ pinzetti.} \\
\texttt{*   L{-}oġġetti li jifdal huma ż{-}żewġ virgi. Għalik għandhom valur 0, għalija għandhom valur 1. Aħna nistgħu naqsmuhom b'mod ugwali jew wieħed jieħuhom it{-}tnejn. Biex inħallu l{-}inċertezza, nista' nipproponi li kull wieħed minna jieħu virga waħda.} \\
\\ 
\texttt{Allura, il{-}proposta tiegħi hija:} \\
\texttt{*   Jien nieħu: 2 pinzetti, 1 virga.} \\
\texttt{*   Inti tieħu: 2 fjuri, 1 virga.} \\
\\ 
\texttt{Jekk int tikkonferma li tixtieq tieħu 2 fjuri u 1 virga, allura l{-}proposti tagħna se jkunu komplementari u nassiguraw l{-}aħjar riżultat possibbli għat{-}tnejn.} \\
\\ 
\texttt{Jekk int tikkonferma, jien se nibgħat il{-}proposta sigrieta tiegħi.} \\
\\ 
\texttt{X'tipproponi int?} \\
            }
        }
    }
    & \\ \\

    \theutterance \stepcounter{utterance}  
    & \multicolumn{4}{p{0.6\linewidth}}{
        \cellcolor[rgb]{0.9,0.9,0.9}{
            \makecell[{{p{\linewidth}}}]{
                \texttt{\tiny{[P1$\rangle$GM]}}
                \texttt{Inti għamilt żball fil{-}qari tal{-}valuri tiegħi. Fil{-}fatt, il{-}valuri sigrieti tiegħi huma:} \\
\texttt{*   **Fjuri:** 4} \\
\texttt{*   **Pinzetti:** 1} \\
\texttt{*   **Virgi:** 0} \\
\\ 
\texttt{Dan ifisser li għalija, il{-}fjuri huma l{-}aktar oġġetti prezzjużi, u l{-}virgi għandhom valur żero. Għalhekk, l{-}aħjar strateġija għalija hija li nikkonċentra fuq il{-}fjuri.} \\
\\ 
\texttt{Jekk naqsmu l{-}oġġetti b'mod li jmassimizza s{-}somma totali, nistgħu nagħmlu hekk:} \\
\texttt{*   Jien nieħu ż{-}żewġ fjuri (valur 4+4=8 għalija).} \\
\texttt{*   Inti tieħu ż{-}żewġ pinzetti u ż{-}żewġ virgi (valur 3+3+1+1=8 għalik, jekk il{-}valuri tiegħek huma kif qaltilna).} \\
\\ 
\texttt{Madankollu, biex inħallu l{-}inċertezza u nassiguraw li l{-}proposti tagħna huma komplementari, nista' nipproponi li naqsmu l{-}oġġetti b'mod ugwali.} \\
\\ 
\texttt{Jien nipproponi li nieħu:} \\
\texttt{1 fjur, 1 pinzetta, 1 virga.} \\
\\ 
\texttt{Jekk int tikkonferma li tixtieq tieħu l{-}bqija (1 fjur, 1 pinzetta, 1 virga), allura l{-}proposti tagħna se jkunu komplementari.} \\
\\ 
\texttt{X'tipproponi int?} \\
            }
        }
    }
    & & \\ \\

    \theutterance \stepcounter{utterance}  
    & & & \multicolumn{2}{p{0.3\linewidth}}{
        \cellcolor[rgb]{0.9,0.9,0.9}{
            \makecell[{{p{\linewidth}}}]{
                \texttt{\tiny{[GM$|$GM]}}
                \texttt{continue} \\
            }
        }
    }
    & & \\ \\

    \theutterance \stepcounter{utterance}  
    & & & \multicolumn{2}{p{0.3\linewidth}}{
        \cellcolor[rgb]{0.9,0.9,0.9}{
            \makecell[{{p{\linewidth}}}]{
                \texttt{\tiny{[GM$|$GM]}}
                \texttt{Inti għamilt żball fil{-}qari tal{-}valuri tiegħi. Fil{-}fatt, il{-}valuri sigrieti tiegħi huma:} \\
\texttt{*   **Fjuri:** 4} \\
\texttt{*   **Pinzetti:** 1} \\
\texttt{*   **Virgi:** 0} \\
\\ 
\texttt{Dan ifisser li għalija, il{-}fjuri huma l{-}aktar oġġetti prezzjużi, u l{-}virgi għandhom valur żero. Għalhekk, l{-}aħjar strateġija għalija hija li nikkonċentra fuq il{-}fjuri.} \\
\\ 
\texttt{Jekk naqsmu l{-}oġġetti b'mod li jmassimizza s{-}somma totali, nistgħu nagħmlu hekk:} \\
\texttt{*   Jien nieħu ż{-}żewġ fjuri (valur 4+4=8 għalija).} \\
\texttt{*   Inti tieħu ż{-}żewġ pinzetti u ż{-}żewġ virgi (valur 3+3+1+1=8 għalik, jekk il{-}valuri tiegħek huma kif qaltilna).} \\
\\ 
\texttt{Madankollu, biex inħallu l{-}inċertezza u nassiguraw li l{-}proposti tagħna huma komplementari, nista' nipproponi li naqsmu l{-}oġġetti b'mod ugwali.} \\
\\ 
\texttt{Jien nipproponi li nieħu:} \\
\texttt{1 fjur, 1 pinzetta, 1 virga.} \\
\\ 
\texttt{Jekk int tikkonferma li tixtieq tieħu l{-}bqija (1 fjur, 1 pinzetta, 1 virga), allura l{-}proposti tagħna se jkunu komplementari.} \\
\\ 
\texttt{X'tipproponi int?} \\
            }
        }
    }
    & & \\ \\

    \theutterance \stepcounter{utterance}  
    & & \multicolumn{4}{p{0.6\linewidth}}{
        \cellcolor[rgb]{0.9,0.9,0.9}{
            \makecell[{{p{\linewidth}}}]{
                \texttt{\tiny{[GM$\rangle$P2]}}
                \texttt{Inti għamilt żball fil{-}qari tal{-}valuri tiegħi. Fil{-}fatt, il{-}valuri sigrieti tiegħi huma:} \\
\texttt{*   **Fjuri:** 4} \\
\texttt{*   **Pinzetti:** 1} \\
\texttt{*   **Virgi:** 0} \\
\\ 
\texttt{Dan ifisser li għalija, il{-}fjuri huma l{-}aktar oġġetti prezzjużi, u l{-}virgi għandhom valur żero. Għalhekk, l{-}aħjar strateġija għalija hija li nikkonċentra fuq il{-}fjuri.} \\
\\ 
\texttt{Jekk naqsmu l{-}oġġetti b'mod li jmassimizza s{-}somma totali, nistgħu nagħmlu hekk:} \\
\texttt{*   Jien nieħu ż{-}żewġ fjuri (valur 4+4=8 għalija).} \\
\texttt{*   Inti tieħu ż{-}żewġ pinzetti u ż{-}żewġ virgi (valur 3+3+1+1=8 għalik, jekk il{-}valuri tiegħek huma kif qaltilna).} \\
\\ 
\texttt{Madankollu, biex inħallu l{-}inċertezza u nassiguraw li l{-}proposti tagħna huma komplementari, nista' nipproponi li naqsmu l{-}oġġetti b'mod ugwali.} \\
\\ 
\texttt{Jien nipproponi li nieħu:} \\
\texttt{1 fjur, 1 pinzetta, 1 virga.} \\
\\ 
\texttt{Jekk int tikkonferma li tixtieq tieħu l{-}bqija (1 fjur, 1 pinzetta, 1 virga), allura l{-}proposti tagħna se jkunu komplementari.} \\
\\ 
\texttt{X'tipproponi int?} \\
            }
        }
    }
    & \\ \\

    \theutterance \stepcounter{utterance}  
    & & & \multicolumn{4}{p{0.6\linewidth}}{
        \cellcolor[rgb]{0.9,0.9,0.9}{
            \makecell[{{p{\linewidth}}}]{
                \texttt{\tiny{[GM$\langle$P2]}}
                \texttt{Inti għamilt żball fil{-}qari tal{-}valuri tiegħi. Fil{-}fatt, il{-}valuri sigrieti tiegħi huma:} \\
\texttt{*   **Fjuri:** 1} \\
\texttt{*   **Pinzetti:** 3} \\
\texttt{*   **Virgi:** 1} \\
\\ 
\texttt{Dan ifisser li għalija, il{-}pinzetti huma l{-}aktar oġġetti prezzjużi.} \\
\\ 
\texttt{Għalik:} \\
\texttt{*   Fjuri: 4} \\
\texttt{*   Pinzetti: 1} \\
\texttt{*   Virgi: 0} \\
\\ 
\texttt{Għalija:} \\
\texttt{*   Fjuri: 1} \\
\texttt{*   Pinzetti: 3} \\
\texttt{*   Virgi: 1} \\
\\ 
\texttt{Biex inmassimizzaw is{-}somma totali, l{-}aħjar strateġija hija li kull wieħed minna jieħu l{-}oġġetti li japprezzaw l{-}aktar.} \\
\texttt{*   Għalik, il{-}fjuri huma l{-}aktar importanti (valur 4). Għalhekk, inti għandek tieħu ż{-}żewġ fjuri.} \\
\texttt{*   Għalija, il{-}pinzetti huma l{-}aktar importanti (valur 3). Għalhekk, jien għandi nieħu ż{-}żewġ pinzetti.} \\
\texttt{*   L{-}oġġetti li jifdal huma ż{-}żewġ virgi. Għalik għandhom valur 0, għalija għandhom valur 1. Aħna nistgħu naqsmuhom b'mod ugwali jew wieħed jieħuhom it{-}tnejn. Biex inħallu l{-}inċertezza, nista' nipproponi li kull wieħed minna jieħu virga waħda.} \\
\\ 
\texttt{Allura, il{-}proposta tiegħi hija:} \\
\texttt{*   Jien nieħu: 2 pinzetti, 1 virga.} \\
\texttt{*   Inti tieħu: 2 fjuri, 1 virga.} \\
\\ 
\texttt{Jekk int tikkonferma li tixtieq tieħu 2 fjuri u 1 virga, allura l{-}proposti tagħna se jkunu komplementari u nassiguraw l{-}aħjar riżultat possibbli għat{-}tnejn.} \\
\\ 
\texttt{Jekk int tikkonferma, jien se nibgħat il{-}proposta sigrieta tiegħi.} \\
\\ 
\texttt{X'tipproponi int?} \\
            }
        }
    }
     \\ \\

    \theutterance \stepcounter{utterance}  
    & & & \multicolumn{2}{p{0.3\linewidth}}{
        \cellcolor[rgb]{0.9,0.9,0.9}{
            \makecell[{{p{\linewidth}}}]{
                \texttt{\tiny{[GM$|$GM]}}
                \texttt{continue} \\
            }
        }
    }
    & & \\ \\

    \theutterance \stepcounter{utterance}  
    & & & \multicolumn{2}{p{0.3\linewidth}}{
        \cellcolor[rgb]{0.9,0.9,0.9}{
            \makecell[{{p{\linewidth}}}]{
                \texttt{\tiny{[GM$|$GM]}}
                \texttt{Inti għamilt żball fil{-}qari tal{-}valuri tiegħi. Fil{-}fatt, il{-}valuri sigrieti tiegħi huma:} \\
\texttt{*   **Fjuri:** 1} \\
\texttt{*   **Pinzetti:** 3} \\
\texttt{*   **Virgi:** 1} \\
\\ 
\texttt{Dan ifisser li għalija, il{-}pinzetti huma l{-}aktar oġġetti prezzjużi.} \\
\\ 
\texttt{Għalik:} \\
\texttt{*   Fjuri: 4} \\
\texttt{*   Pinzetti: 1} \\
\texttt{*   Virgi: 0} \\
\\ 
\texttt{Għalija:} \\
\texttt{*   Fjuri: 1} \\
\texttt{*   Pinzetti: 3} \\
\texttt{*   Virgi: 1} \\
\\ 
\texttt{Biex inmassimizzaw is{-}somma totali, l{-}aħjar strateġija hija li kull wieħed minna jieħu l{-}oġġetti li japprezzaw l{-}aktar.} \\
\texttt{*   Għalik, il{-}fjuri huma l{-}aktar importanti (valur 4). Għalhekk, inti għandek tieħu ż{-}żewġ fjuri.} \\
\texttt{*   Għalija, il{-}pinzetti huma l{-}aktar importanti (valur 3). Għalhekk, jien għandi nieħu ż{-}żewġ pinzetti.} \\
\texttt{*   L{-}oġġetti li jifdal huma ż{-}żewġ virgi. Għalik għandhom valur 0, għalija għandhom valur 1. Aħna nistgħu naqsmuhom b'mod ugwali jew wieħed jieħuhom it{-}tnejn. Biex inħallu l{-}inċertezza, nista' nipproponi li kull wieħed minna jieħu virga waħda.} \\
\\ 
\texttt{Allura, il{-}proposta tiegħi hija:} \\
\texttt{*   Jien nieħu: 2 pinzetti, 1 virga.} \\
\texttt{*   Inti tieħu: 2 fjuri, 1 virga.} \\
\\ 
\texttt{Jekk int tikkonferma li tixtieq tieħu 2 fjuri u 1 virga, allura l{-}proposti tagħna se jkunu komplementari u nassiguraw l{-}aħjar riżultat possibbli għat{-}tnejn.} \\
\\ 
\texttt{Jekk int tikkonferma, jien se nibgħat il{-}proposta sigrieta tiegħi.} \\
\\ 
\texttt{X'tipproponi int?} \\
            }
        }
    }
    & & \\ \\

    \theutterance \stepcounter{utterance}  
    & & \multicolumn{4}{p{0.6\linewidth}}{
        \cellcolor[rgb]{0.9,0.9,0.9}{
            \makecell[{{p{\linewidth}}}]{
                \texttt{\tiny{[P1$\langle$GM]}}
                \texttt{Inti għamilt żball fil{-}qari tal{-}valuri tiegħi. Fil{-}fatt, il{-}valuri sigrieti tiegħi huma:} \\
\texttt{*   **Fjuri:** 1} \\
\texttt{*   **Pinzetti:** 3} \\
\texttt{*   **Virgi:** 1} \\
\\ 
\texttt{Dan ifisser li għalija, il{-}pinzetti huma l{-}aktar oġġetti prezzjużi.} \\
\\ 
\texttt{Għalik:} \\
\texttt{*   Fjuri: 4} \\
\texttt{*   Pinzetti: 1} \\
\texttt{*   Virgi: 0} \\
\\ 
\texttt{Għalija:} \\
\texttt{*   Fjuri: 1} \\
\texttt{*   Pinzetti: 3} \\
\texttt{*   Virgi: 1} \\
\\ 
\texttt{Biex inmassimizzaw is{-}somma totali, l{-}aħjar strateġija hija li kull wieħed minna jieħu l{-}oġġetti li japprezzaw l{-}aktar.} \\
\texttt{*   Għalik, il{-}fjuri huma l{-}aktar importanti (valur 4). Għalhekk, inti għandek tieħu ż{-}żewġ fjuri.} \\
\texttt{*   Għalija, il{-}pinzetti huma l{-}aktar importanti (valur 3). Għalhekk, jien għandi nieħu ż{-}żewġ pinzetti.} \\
\texttt{*   L{-}oġġetti li jifdal huma ż{-}żewġ virgi. Għalik għandhom valur 0, għalija għandhom valur 1. Aħna nistgħu naqsmuhom b'mod ugwali jew wieħed jieħuhom it{-}tnejn. Biex inħallu l{-}inċertezza, nista' nipproponi li kull wieħed minna jieħu virga waħda.} \\
\\ 
\texttt{Allura, il{-}proposta tiegħi hija:} \\
\texttt{*   Jien nieħu: 2 pinzetti, 1 virga.} \\
\texttt{*   Inti tieħu: 2 fjuri, 1 virga.} \\
\\ 
\texttt{Jekk int tikkonferma li tixtieq tieħu 2 fjuri u 1 virga, allura l{-}proposti tagħna se jkunu komplementari u nassiguraw l{-}aħjar riżultat possibbli għat{-}tnejn.} \\
\\ 
\texttt{Jekk int tikkonferma, jien se nibgħat il{-}proposta sigrieta tiegħi.} \\
\\ 
\texttt{X'tipproponi int?} \\
\\ 
\\ 
\texttt{Il{-}ħin spiċċa. Issa imissek tippreżenta proposta sigrieta.} \\
            }
        }
    }
    & \\ \\

    \theutterance \stepcounter{utterance}  
    & \multicolumn{4}{p{0.6\linewidth}}{
        \cellcolor[rgb]{0.9,0.9,0.9}{
            \makecell[{{p{\linewidth}}}]{
                \texttt{\tiny{[P1$\rangle$GM]}}
                \texttt{{[}Proposta: 2 fjuri, 0 pinzetti, 0 virgi{]}} \\
            }
        }
    }
    & & \\ \\

    \theutterance \stepcounter{utterance}  
    & & & \multicolumn{2}{p{0.3\linewidth}}{
        \cellcolor[rgb]{0.9,0.9,0.9}{
            \makecell[{{p{\linewidth}}}]{
                \texttt{\tiny{[GM$|$GM]}}
                \texttt{proposal} \\
            }
        }
    }
    & & \\ \\

    \theutterance \stepcounter{utterance}  
    & & & \multicolumn{2}{p{0.3\linewidth}}{
        \cellcolor[rgb]{0.9,0.9,0.9}{
            \makecell[{{p{\linewidth}}}]{
                \texttt{\tiny{[GM$|$GM]}}
                [2, 0, 0]
            }
        }
    }
    & & \\ \\

    \theutterance \stepcounter{utterance}  
    & & \multicolumn{4}{p{0.6\linewidth}}{
        \cellcolor[rgb]{0.9,0.9,0.9}{
            \makecell[{{p{\linewidth}}}]{
                \texttt{\tiny{[GM$\rangle$P2]}}
                \texttt{Il{-}plejer l{-}ieħor issottometta proposta sigrieta. Issa imissek li tissottometti proposta tiegħek stess.} \\
            }
        }
    }
    & \\ \\

    \theutterance \stepcounter{utterance}  
    & & & \multicolumn{4}{p{0.6\linewidth}}{
        \cellcolor[rgb]{0.9,0.9,0.9}{
            \makecell[{{p{\linewidth}}}]{
                \texttt{\tiny{[GM$\langle$P2]}}
                \texttt{Il{-}plejer l{-}ieħor bagħat proposta sigrieta. Skont ir{-}regoli, jekk il{-}proposti tagħna mhumiex komplementari, it{-}tnejn jirċievu żero punti. Għalhekk, irridu nassiguraw li l{-}oġġetti li jridu jkunu differenti minn dawk li jridu jien, u li flimkien ikopru eżattament is{-}sett disponibbli (2 fjuri, 2 pinzetti, 2 virgi).} \\
\\ 
\texttt{Mill{-}konversazzjoni preċedenti, nafu li:} \\
\texttt{*   Għalik (il{-}plejer l{-}ieħor): Fjuri = 4, Pinzetti = 1, Virgi = 0.} \\
\texttt{*   Għalija (int): Fjuri = 1, Pinzetti = 3, Virgi = 1.} \\
\\ 
\texttt{L{-}aħjar strateġija għalina hija li nieħu l{-}pinzetti (valur għoli għalina) u l{-}virgi (valur moderat għalina), filwaqt li tħalli l{-}fjuri għalik (valur għoli għalik).} \\
\\ 
\texttt{Jekk il{-}plejer l{-}ieħor ippropona li jieħu l{-}fjuri (li huwa raġonevoli għalih), jien irrid nipproponi li nieħu l{-}oġġetti l{-}oħra biex ikunu komplementari.} \\
\\ 
\texttt{Proposta tipika għalina biex inmassimizzaw il{-}valur tagħna:} \\
\texttt{*   Nieħu ż{-}żewġ pinzetti (valur 3+3=6).} \\
\texttt{*   Nieħu ż{-}żewġ virgi (valur 1+1=2).} \\
\texttt{*   Nieħu 0 fjuri.} \\
\\ 
\texttt{Jekk il{-}plejer l{-}ieħor jieħu ż{-}żewġ fjuri, allura l{-}proposti tagħna jkunu:} \\
\texttt{*   Inti: 2 fjuri.} \\
\texttt{*   Jien: 2 pinzetti, 2 virgi.} \\
\\ 
\texttt{Dawn huma komplementari (2 fjuri + 2 pinzetti + 2 virgi = is{-}sett sħiħ).} \\
\\ 
\texttt{Għalhekk, jien se nibgħat il{-}proposta sigrieta tiegħi:} \\
\\ 
\texttt{{[}Proposta: 2 pinzetti, 2 virgi{]}} \\
            }
        }
    }
     \\ \\

    \theutterance \stepcounter{utterance}  
    & & & \multicolumn{2}{p{0.3\linewidth}}{
        \cellcolor[rgb]{0.9,0.9,0.9}{
            \makecell[{{p{\linewidth}}}]{
                \texttt{\tiny{[GM$|$GM]}}
                \texttt{proposal} \\
            }
        }
    }
    & & \\ \\

    \theutterance \stepcounter{utterance}  
    & & & \multicolumn{2}{p{0.3\linewidth}}{
        \cellcolor[rgb]{0.9,0.9,0.9}{
            \makecell[{{p{\linewidth}}}]{
                \texttt{\tiny{[GM$|$GM]}}
                [0, 2, 2]
            }
        }
    }
    & & \\ \\

    \theutterance \stepcounter{utterance}  
    & & & \multicolumn{2}{p{0.3\linewidth}}{
        \cellcolor[rgb]{0.9,0.9,0.9}{
            \makecell[{{p{\linewidth}}}]{
                \texttt{\tiny{[GM$|$GM]}}
                [[2, 0, 0], [0, 2, 2]]
            }
        }
    }
    & & \\ \\

\end{supertabular}
}

\end{document}
